\documentclass[10pt,a4paper]{amsart}
\usepackage{amsmath,amssymb,amsthm,graphicx}
\usepackage{verbatim}
\newenvironment{reference}{\comment}{\endcomment}
%\newenvironment{reference}{}{}
\newenvironment{slogan}{\comment}{\endcomment}
\newenvironment{history}{\comment}{\endcomment}

% For commutative diagrams we use Xy-pic
\usepackage[all]{xy}

% We use 2cell for 2-commutative diagrams.
\xyoption{2cell}
% Theorem environments.
%
\theoremstyle{plain}
\newtheorem{theorem}[subsection]{Theorem}
\newtheorem{proposition}[subsection]{Proposition}
\newtheorem{lemma}[subsection]{Lemma}

\theoremstyle{definition}
\newtheorem{definition}[subsection]{Definition}
\newtheorem{example}[subsection]{Example}
\newtheorem{exercise}[subsection]{Exercise}
\newtheorem{situation}[subsection]{Situation}

\theoremstyle{remark}
\newtheorem{remark}[subsection]{Remark}
\newtheorem{remarks}[subsection]{Remarks}

\numberwithin{equation}{subsection}

% Macros
%
\def\lim{\mathop{\mathrm{lim}}\nolimits}
\def\colim{\mathop{\mathrm{colim}}\nolimits}
\def\Spec{\mathop{\mathrm{Spec}}}
\def\Hom{\mathop{\mathrm{Hom}}\nolimits}
\def\Ext{\mathop{\mathrm{Ext}}\nolimits}
\def\SheafHom{\mathop{\mathcal{H}\!\mathit{om}}\nolimits}
\def\SheafExt{\mathop{\mathcal{E}\!\mathit{xt}}\nolimits}
\def\Sch{\mathit{Sch}}
\def\Mor{\mathop{\mathrm{Mor}}\nolimits}
\def\Ob{\mathop{\mathrm{Ob}}\nolimits}
\def\Sh{\mathop{\mathit{Sh}}\nolimits}
\def\NL{\mathop{N\!L}\nolimits}
\def\CH{\mathop{\mathrm{CH}}\nolimits}
\def\proetale{{pro\text{-}\acute{e}tale}}
\def\etale{{\acute{e}tale}}
\def\QCoh{\mathit{QCoh}}
\def\Ker{\mathop{\mathrm{Ker}}}
\def\Im{\mathop{\mathrm{Im}}}
\def\Coker{\mathop{\mathrm{Coker}}}
\def\Coim{\mathop{\mathrm{Coim}}}

% Boxtimes
%
\DeclareMathSymbol{\boxtimes}{\mathbin}{AMSa}{"02}

%
% Macros for moduli stacks/spaces
%
\def\QCohstack{\mathcal{QC}\!\mathit{oh}}
\def\Cohstack{\mathcal{C}\!\mathit{oh}}
\def\Spacesstack{\mathcal{S}\!\mathit{paces}}
\def\Quotfunctor{\mathrm{Quot}}
\def\Hilbfunctor{\mathrm{Hilb}}
\def\Curvesstack{\mathcal{C}\!\mathit{urves}}
\def\Polarizedstack{\mathcal{P}\!\mathit{olarized}}
\def\Complexesstack{\mathcal{C}\!\mathit{omplexes}}
% \Pic is the operator that assigns to X its picard group, usage \Pic(X)
% \Picardstack_{X/B} denotes the Picard stack of X over B
% \Picardfunctor_{X/B} denotes the Picard functor of X over B
\def\Pic{\mathop{\mathrm{Pic}}\nolimits}
\def\Picardstack{\mathcal{P}\!\mathit{ic}}
\def\Picardfunctor{\mathrm{Pic}}
\def\Deformationcategory{\mathcal{D}\!\mathit{ef}}
\usepackage{cleveref}
\crefname{originallemma}{original lemma}{original lemmas}
\Crefname{originallemma}{Original lemma}{Original lemmas}
\makeatletter
\newlabel{lemma-pushforward}{{03M9}{source}{}{sourcepointer-1}{}}
\newlabel{lemma-pushforward@cref}{{[originallemma][0][]03M9}{[1][0][]source}{}{}{}}
\makeatother
\crefname{originallemma}{original lemma}{original lemmas}
\Crefname{originallemma}{Original lemma}{Original lemmas}
\makeatletter
\newlabel{lemma-scheme-theoretic-image}{{082X}{source}{}{sourcepointer-2}{}}
\newlabel{lemma-scheme-theoretic-image@cref}{{[originallemma][0][]082X}{[1][0][]source}{}{}{}}
\makeatother
\title{The scheme theoretic image of a quasi-compact morphism of algebraic spaces}
\author{The Stacks Project Authors}
\date{stacks-project:spaces-morphisms.tex:local-snapshot-20261012; GNU FDL 1.2 or later}
\begin{document}
\maketitle
\setcounter{subsection}{1}
\noindent\textit{Editorial note.} All delivered mathematical text is the author's own, copied from the pinned Stacks Project chapter file spaces-morphisms.tex; no statement or proof step was rewritten, reordered or removed.\par\medskip
\noindent\textit{Editorial note.} The author refers back to a computation performed in the proof of the scheme theoretic image lemma for assertion (2); that reference is to a statement in the same chapter and the counted proof supplies the quasi-coherence argument for I in full.\par\medskip
\noindent\textit{Editorial note.} The two references of the counted spans are recorded as non-load-bearing navigation pointers; the argument that consumes them stays inside the counted proof span.\par\medskip
\noindent\textit{Editorial note.} Preamble: the three native excerpt spans of sources/stacks-spaces-morphisms/preamble.tex (lines 8-13, 14-18 and 154-217). The chapter file itself carries no preamble text.\par\medskip
\noindent\textit{Allowed standard prerequisites (editorial).}\begin{itemize}
\item The scheme theoretic image of a quasi-compact morphism and its defining ideal. Lemma lemma-scheme-theoretic-image, used at proof lines 3268-3271. Supplies the identification of Z with the closed subspace defined by I once (1) holds.
\item The pushforward of the structure sheaf along a quasi-compact quasi-separated morphism is quasi-coherent. Lemma lemma-pushforward, used at proof lines 3276-3280. Makes f'\_*O\_U quasi-coherent, so that its kernel is.
\item Affine etale covers of quasi-compact morphisms and etale locality of quasi-coherence. used at proof lines 3272-3275. Reduces quasi-coherence of I to the case of an affine base.
\item Density of the image in the scheme theoretic image. used at proof lines 3281-3283. Derives assertion (4) from the description of I.
\end{itemize}\medskip
\ifdefined\hypertarget\hypertarget{sourcepointer-1}{}\fi
\noindent\textit{Source reference.} Original lemma 03M9: focused-seeds/AGP/sources/stacks-snapshot-20261011/tags.txt, source lines 4741--4741. Navigation pointer into the Stacks Project; the standard fact it supplies is applied inside the counted proof span and the argument resting on it stays there.\par\medskip
\ifdefined\hypertarget\hypertarget{sourcepointer-2}{}\fi
\noindent\textit{Source reference.} Original lemma 082X: focused-seeds/AGP/sources/stacks-snapshot-20261011/tags.txt, source lines 10267--10267. Navigation pointer into the Stacks Project; the standard fact it supplies is applied inside the counted proof span and the argument resting on it stays there.\par\medskip
\begin{lemma}
\label{lemma-quasi-compact-scheme-theoretic-image}
Let $S$ be a scheme.
Let $f : X \to Y$ be a morphism of algebraic spaces over $S$.
Let $Z \subset Y$ be the scheme theoretic image of $f$.
If $f$ is quasi-compact then
\begin{enumerate}
\item the sheaf of ideals
$\mathcal{I} = \Ker(\mathcal{O}_Y \to f_*\mathcal{O}_X)$
is quasi-coherent,
\item the scheme theoretic image $Z$ is the closed subspace
corresponding to $\mathcal{I}$,
\item for any \'etale morphism $V \to Y$ the scheme theoretic image of
$X \times_Y V \to V$ is equal to $Z \times_Y V$, and
\item the image $|f|(|X|) \subset |Z|$ is a dense subset of $|Z|$.
\end{enumerate}
\end{lemma}

\begin{proof}
To prove (3) it suffices to prove (1) and (2) since the
formation of $\mathcal{I}$ commutes with \'etale localization.
If (1) holds then in the proof of Lemma \ref{lemma-scheme-theoretic-image}
we showed (2). Let us prove that $\mathcal{I}$ is quasi-coherent.
Since the property of being quasi-coherent is \'etale local we may
assume $Y$ is an affine scheme. As $f$ is quasi-compact,
we can find an affine scheme $U$ and a surjective \'etale morphism
$U \to X$. Denote $f'$ the composition $U \to X \to Y$.
Then $f_*\mathcal{O}_X$ is a subsheaf of $f'_*\mathcal{O}_U$,
and hence $\mathcal{I} = \Ker(\mathcal{O}_Y \to \mathcal{O}_{X'})$.
By Lemma \ref{lemma-pushforward}
the sheaf $f'_*\mathcal{O}_U$ is quasi-coherent on $Y$. Hence $\mathcal{I}$
is quasi-coherent as a kernel of a map between coherent modules.
Finally, part (4) follows from parts (1), (2), and (3) as the ideal
$\mathcal{I}$ will be the unit ideal in any point of $|Y|$ which is
not contained in the closure of $|f|(|X|)$.
\end{proof}

\end{document}
